% Generated draft: requires mathematical and source review.
\documentclass[11pt]{article}
\usepackage[margin=25mm]{geometry}
\usepackage{fontspec}
\setmainfont{Latin Modern Roman}
\usepackage{amsmath,amssymb,mathtools,unicode-math}
\setmathfont{Latin Modern Math}
\usepackage{enumitem,needspace,xurl,hyperref}
\setlength{\parindent}{0pt}
\setlength{\parskip}{5pt}
\setlength{\emergencystretch}{3em}
\newcommand{\N}{\mathbb{N}}
\newcommand{\Z}{\mathbb{Z}}
\newcommand{\Q}{\mathbb{Q}}
\newcommand{\R}{\mathbb{R}}
\newcommand{\C}{\mathbb{C}}
\newcommand{\F}{\mathbb{F}}
\newcommand{\A}{\mathbb{A}}
\newcommand{\PP}{\mathbb{P}}
\newcommand{\E}{\mathbb{E}}
\newcommand{\eps}{\varepsilon}
\newcommand{\dd}{\,\mathrm{d}}
\begin{document}
\section*{UnsolvedMath ID: 20001145}
\textbf{Source problem code:} catalogue entry
\subsection{Definitions and mathematical statement}
This is an aggregate source record, not a single conjecture. It starts inside a table of contents, includes substantive text from Chapters A--E, and ends immediately after the definition in E.8. The supplied verification identifies the document as \emph{Perfect Graphs}, version 24 August 2004. The following preserves the actual questions under their source labels; it does not supply the missing later chapters or the questions following E.8.

Let $G=(V,E)$ be a finite simple graph, $\overline G$ its complement, $G[S]$ its induced subgraph, $\omega(G)$ its maximum clique size, and $\chi(G)$ its chromatic number. A graph is perfect if $\chi(G[S])=\omega(G[S])$ for every $S\subseteq V$. A hole is an induced cycle of length at least four, an antihole the complement of such a cycle, and a Berge graph has no odd hole or odd antihole. A clique and a proper coloring with the same number of colors certify equality for that graph; this certificate for $G$ alone does not certify all its induced subgraphs. An even pair consists of two distinct nonadjacent vertices for which every induced path joining them has even length.

A skew partition in the four-set notation of E.5 is a partition $(A,B,C,D)$ into nonempty sets with $A$ anticomplete to $B$ and $C$ complete to $D$. The E.8 convention exchanges these roles: $A$ is complete to $B$, and $C$ is anticomplete to $D$. A star cutset is a vertex cutset contained in the closed neighborhood of one of its members. Boundary conventions for disconnected graphs and empty disk structures are not specified in the source.

\begin{enumerate}
\item \textbf{A.1, A.4, A.5, A.5.a: recognition and certificates.} The chapter asks for polynomial-time perfect-graph recognition. A.1 reports that Berge recognition was obtained in November 2002, independently of the proof of the Strong Perfect Graph Theorem, by the two groups named in the sources. A.4 (Jack Edmonds) asks for an NP description of perfect graphs, that is, polynomially verifiable polynomial-length certificates of perfectness. The supplied research aid says this literal request is already satisfied by polynomial Berge recognition, reported there as $O(|V|^9)$, together with the Strong Perfect Graph Theorem. This is an attributed status qualification, not a new literature assessment. A.5 (Bruce Shepherd) asks for polynomial recognition given the list of all maximal cliques and explicitly reports this resolved by general recognition. The vertex set, including isolated vertices, and completeness of the supplied list must be known. A.5.a (Jeremy Spinrad) asks for polynomial recognition of Berge graphs having polynomially bounded numbers of maximal cliques. These restricted requests are retained despite the general recognition result.
\item \textbf{A.2: interaction of decompositions.} Jeremy Spinrad asks whether odd-hole testing using mixed skew-partition and $2$-join decompositions can be reduced to tests using only skew partitions and tests using only $2$-joins. Kristina Vu\v{s}kovi\'c describes the analogous even-hole-free strategy: first decompose by vertex cutsets, then by $2$-joins whose blocks introduce no new vertex cutset. Neither a precise reduction nor its permitted blocks are specified.
\item \textbf{A.3: robust optimization.} Kathie Cameron and Jack Edmonds ask for a combinatorial polynomial-time algorithm which, for every graph $G$, returns either a clique $K$ and proper coloring $c$ with $|K|=|c(V)|$, or an easily checked imperfection obstruction. Two alternatives for the obstruction are an ``$\alpha$--$\omega$ partitioned subgraph'' (not defined in this extraction), and an odd hole or odd antihole. The two stated existential alternatives allow both outputs to exist; in particular, a graph that is not perfect may still admit an equal-size clique and coloring. This task therefore need not recognize every imperfect input as imperfect.

The source motivates this with Edmonds' robust Meyniel algorithm: return an odd cycle of length greater than three with at most one chord, or an equal-size clique and coloring. It contrasts the algorithms of Hoang and Hertz, which assume the Meyniel property, and Burlet--Fonlupt decomposition recognition. Its matrix analogy asks, for an arbitrary $0$--$1$ matrix $M$ with maximum row sum $x$, for either an $x$-coloring of columns making all ones in each row differently colored, or an odd-hole submatrix. Here the usual matrix obstruction is a square submatrix of odd order at least three with exactly two ones in every row and column. The source contrasts Cameron--Edmonds' robust method with Conforti--Cornu\'ejols recognition and Conforti--Cornu\'ejols--Rao decomposition, and suggests the analogous relation to Chudnovsky--Robertson--Seymour--Thomas Berge decomposition.
\item \textbf{A.6: TDI and local unimodularity.} Given $A\in\{0,1\}^{m\times n}$ and a vector $b$, decide whether $Ax\leq b$ is totally dual integral: for every integral objective $c\in\Z^n$ for which the dual minimization $\min\{b^Ty:A^Ty=c,\ y\geq0\}$ has a finite attained optimum, an integral optimum $y$ exists. The source then considers the separate polyhedron $P=\{x:Ax\leq b,\ x\geq0\}$ and, for each vertex $u$, the matrix $A(u)$ of primitive integral facet normals of the facets containing $u$. Its positive-$b$ input question is whether every $A(u)$ is square of determinant one, and whether this can be decided in polynomial time. An ordering-independent unimodularity condition would be $|\det A(u)|=1$; the literal determinant sign is ambiguous. The source attributes a corresponding assertion about ``nondegenerate'' TDI matrices to Serkan Ho\c{s}ten, without defining nondegeneracy or the representation hypotheses needed for that assertion. Andr\'as Seb\H{o} suggests reducing matrix perfectness by lexicographically perturbing the all-ones vector to obtain $b$. The extraction calls that vector an objective function although $b$ is a right-hand side. No particular perturbation, encoding of $b$, or equivalence between the two displayed systems is supplied.
\item \textbf{A.7: parameterized algorithm.} Mohammad Taghi Hajiaghayi asks for an algorithm of running time $O(f(k)|V|^c)$, where $k=\omega(G)$, $c$ is a small constant independent of $k$, and $f$ is described as exponential. The exact task to be solved, beyond its perfect-graph algorithmic context, is not restated. Downey and Fellows are credited for the parameterized-algorithm framework.
\item \textbf{A.8: clique joins.} Put $\omega=\omega(G)$. A $k$-clique-join consists of $(A_0,B_0),(A_1,B_1),\ldots,(A_k,B_k)$, where $A_0,B_0$ partition $V$, each contains an $\omega$-clique, $A_i\subseteq A_0$, $B_i\subseteq B_0$, $A_i$ is complete to $B_i$, and every $\omega$-clique meeting both sides is contained in some $A_i\cup B_i$. Overlap among the $A_i$ and among the $B_i$ is allowed. Andr\'as Seb\H{o} asks whether the minimum possible $k$ can be computed or well characterized, both generally and for Berge graphs, and whether a variant gives a composition operation preserving perfectness. The source reports that partitionable graphs have no such join for $k<2(\omega-1)$ and that odd holes and odd antiholes have joins with $k=2(\omega-1)$. ``Partitionable'' and the convention when no join exists are not specified here.
\item \textbf{A.9 and B: structural decomposition.} Bruce Reed asks for new polynomially recognizable endblocks yielding polynomial-size skew-partition or star-cutset decomposition trees. Kristina Vu\v{s}kovi\'c asks whether using only two of the usual four skew-partition blocks suffices for a polynomial tree, possibly with new endblocks or cleaning, and whether star-cutset decomposition of $C_4$-free graphs yields a polynomial tree, possibly with cleaning. The choice of blocks and cleaning operations is not specified. Chapter B (Paul Seymour) asks for explicit structural constructions for subclasses of Berge graphs.
\item \textbf{C and C.1: coloring and unique colorability.} The chapter asks for efficient coloring of perfect graphs. A perfect graph is uniquely colorable if its partition into $\omega(G)$ nonempty stable sets is unique up to permuting the parts. A minimally imperfect graph is imperfect with every proper induced subgraph perfect. The source cites Padberg for the statement that $G-v$ is uniquely colorable for every vertex of a minimally imperfect graph, and describes unique-colorability recognition on perfect graphs by an ellipsoid algorithm in the cited IPCO volume.

Jean Fonlupt and Andr\'as Seb\H{o} ask whether a uniquely colorable perfect graph has two distinct $\omega$-cliques intersecting in $\omega-1$ vertices. The two vertices in their symmetric difference are called forced. They ask whether every known uniquely colorable perfect graph collapses to an $\omega$-clique by successive identifications of forced vertices. The phrase ``every known'' does not define a quantified graph class, and the first question lacks a noncomplete-graph exception. The text also mentions, without describing the positions, that three suitably positioned forced vertices in a minimally imperfect graph imply an odd hole or antihole. Its stated further target is: every minimally imperfect graph $G$ has a vertex $v$ for which $G[N_G(v)]$ is uniquely colorable; it says a simpler proof would shortcut the proof of the Strong Perfect Graph Theorem. The omitted positions are not reconstructed.
\item \textbf{D and D.1: optimization.} The chapter seeks optimization on perfect graphs without the ellipsoid method. Mohammad Hajiaghayi asks for further optimization problems, beyond coloring and maximum clique size, that are easier on perfect graphs, and whether existing or future recognition algorithms help. D.1.a (Bruce Reed) asks, for a perfect graph and two specified vertices, whether there is a hole containing both vertices.
\item \textbf{E.1--E.4: skew partitions and star cutsets.} Paul Seymour asks for a structural obstruction theorem characterizing when a given skew partition of an induced subgraph extends to one of $G$; the source attributes an algorithmic answer to de Figueiredo, Klein, Kohayakawa and Reed. He asks for constructions or a structure theorem for graphs with no skew partition, perhaps by additions preserving that property, and whether finding skew partitions is easier for Berge graphs. Bruce Reed asks the analogous structural construction question for graphs without star cutsets.

Ryan Hayward's stated conjecture is: if neither $G$ nor $\overline G$ has a star cutset, then the disk structure of $G$ is connected. Its vertices are all holes and antiholes of $G$, with two disks adjacent when they share at least two vertices. The source leaves the empty-disk case and possible small-graph restrictions unstated.
\item \textbf{E.5: compatible skew partitions.} Given $t$ skew partitions $(A_i,B_i,C_i,D_i)$ in the anticomplete-$A_i,B_i$ convention, choose one part $X_i$ from each and form the chunk $G[V\setminus\bigcup_i X_i]$. The source states that any odd hole or odd antihole is contained in some chunk. It asks for a compatibility condition, analogous to noncrossing separations in excluded-minor decompositions, producing only linearly or polynomially many necessary chunks rather than up to $4^t$, even when $t=O(|V|)$. The required compatibility condition is the unknown target.
\item \textbf{E.6: size balance.} For a fixed constant $c>0$, call a skew partition size-balanced when each of its four parts has at least $|V|/c$ vertices. Jeremy Spinrad asks (a) whether such a partition can be found in polynomial time whenever one exists, and (b) whether always maximizing the smallest part produces a polynomial-size recursive decomposition tree for perfect graphs. The quantification of $c$, the precise recursion, and the handling of graphs with no skew partition are not fully specified. The source gives as an obstruction family a clique with additional degree-two vertices, each adjacent precisely to the endpoints of one selected clique edge, and says this shows that size-balanced partitions need not exist.

The supplied research aid isolates a structural certificate for this family: all added vertices followed by the original clique form a perfect elimination ordering, meaning each vertex's later neighbors form a clique. Such an ordering certifies perfection of every induced subgraph by restriction. For a clique of size $m\geq3$, the same family admits a clique and coloring of size $m$; if $m=2$ and at least one added vertex is present, the optimum is three. This is a special-family certificate, not a proposed perfect-elimination certificate for all perfect graphs, nor a resolution of the recursive size-balance questions.
\item \textbf{E.7 and E.8: two other balance notions.} E.7 (Jeremy Spinrad) asks whether a supplied skew partition can be checked for being ``balanced'' in polynomial time, but gives no definition here. In particular, its relationship to E.6 size balance is not specified. E.8 defines an even-pair skew partition as a four-set partition $(A,B,C,D)$ with $A$ complete to $B$, $C$ anticomplete to $D$, and every two nonadjacent vertices lying both in $A$ or both in $B$ forming an even pair in $G$. The record ends at this definition; no subsequent E.8 question is supplied.
\end{enumerate}

The opening contents fragment also names Chapters F--L, but supplies no substantive statements from them: F, even pairs in Berge graphs (coloring using even pairs, recognition of even pairs, quasi-parity and strict quasi-parity including forbidden subgraphs and recognition, perfectly contractile graphs and decomposition, possible Berge structure, odd holes and odd walks); G, forbidding holes and antiholes ($2$-divisible graphs, clique coloring, recognition of odd-hole-free graphs, even-hole-free graphs and circulants, and $\beta$-perfect graphs); H, partitionable graphs (perfect/partitionable/kernel-solvable graphs, odd holes, a property, and small transversals); I, imperfection ratio; J, integer programming (partitionable graphs as cutting planes for packing, and theta-body feasibility/membership); K, balanced graphs and balanced circulants; L, $P_4$-structure and relatives. These headings are retained as extraction evidence, not converted into invented problem statements.
\subsection{Short English statement}
This malformed workshop extract contains many separate questions about recognizing, coloring, decomposing, and optimizing perfect graphs. Each substantive question is retained with its source label and attribution. Several recognition requests were already described as resolved; other requests lack definitions or hypotheses, and the final even-pair passage contains only a definition. The supplied machine research singles out the literal NP-certificate request and a special chordal family, but does not repair the aggregate into one mathematical problem.
\subsection{Sources}
\begin{itemize}
\item Local UnsolvedMath registry, permanent ID 20001145, exact source code catalogue entry; canonical collection \texttt{aim-combinatorics-notes.json}, \texttt{notes[269]}, source item 4, tag ``problem'', workshop \emph{The Perfect Graph Conjecture}. Supplied document: \emph{Perfect Graphs}, 21 pages, version 24 August 2004; \url{https://aimath.org/WWN/perfectgraph/perfectgraph.pdf}. No external source was visited for this draft.
\item A.1 credits the November 2002 recognition work to Chudnovsky--Seymour and Cornu\'ejols--Liu--Vu\v{s}kovi\'c, with a paper then due in \emph{Combinatorica}. A.3 names Chudnovsky--Robertson--Seymour--Thomas for Berge decomposition and the Strong Perfect Graph Theorem.
\item A.3 cites Cameron and Edmonds by the supplied identifier [92i:68043], and names Edmonds, Hoang, Hertz, Burlet--Fonlupt, Conforti--Cornu\'ejols, Cameron--Edmonds, and Conforti--Cornu\'ejols--Rao in the robust graph/matrix algorithm comparisons. No full bibliographic entries for those items are supplied.
\item C.1 cites Padberg, ``Perfect zero-one matrices'', \emph{Mathematical Programming} 6 (1974), and the unique-colorability treatment in \emph{IPCO 1}, Kannan and Pulleyblank, editors, Waterloo University Press, 1990; no chapter author or page range is supplied. E.1 cites de Figueiredo, Klein, Kohayakawa and Reed by [2001j:05114]. A.7 credits Downey and Fellows; A.6 attributes its nondegeneracy observation to Serkan Ho\c{s}ten.
\item Individual question contributors, as supplied: Jeremy Spinrad, Kristina Vu\v{s}kovi\'c, Kathie Cameron, Jack Edmonds, Bruce Shepherd, Andr\'as Seb\H{o}, Mohammad Taghi Hajiaghayi (also called Mohammad Hajiaghayi), Bruce Reed, Paul Seymour, Jean Fonlupt, and Ryan Hayward. The subproblem attribution is retained in the mathematical text; spelling normalization does not add bibliographic information.
\item Local research summary and verification: registry label ``open'', research classification OPEN-TRIAGE, statement status ``unrecoverable''; attempt 1, result ``invalid\_statement'', problem status at run ``not\_a\_problem'', recorded rigor ``proved'', novelty confidence ``low'', candidate type ``certificate\_synthesis''. These are supplied metadata, not certification by this draft. The source explicitly warns that the machine-generated aid is not peer reviewed and must be independently verified before citation. Original remarks are empty and the literature field is empty.
\end{itemize}
\end{document}
